\section*{Лекция 2. Функции потерь. Регуляризация}
\subsection* {Введение}

Ранее мы уже научились представлять объекты в данных как некоторые
вектора. После этого ввели модель и показали, что модель линейной регрессии $a$ можно описать как скалярное произведение вектора весов на вектор  признаков:
\[
a: \mathbb{X} \rightarrow \mathbb{Y}
\]

\[
a(x) = \langle w,x\rangle
\]
(при условии, что первый признак "--- единичный).

Мы рассмотрели самые основы, не вдаваясь в более сложную систему, которая
за всем этим стоит. Естественно, мы изучили далеко не все аспекты такого сложного процесса, как \textit{обучения} модели. В частности, непонятными остаются следующие пункты:
\begin{itemize}
    \item Неужели подбор весов происходит перебором? Можно ли сделать это
    как"=то оптимальнее?
    \item В какой момент мы понимаем, что веса оптимальные?
    \item Как мы именно считаем отклонение предсказаний от реальных
    ответов?
    \item \dots
\end{itemize}

Целью наших последующих занятий является постепенный ответ на эти вопросы,
тем самым собирая более полную картину происходящего и постепенно снимая занавес с этого таинственного процесса обучения модели.